\chapter{Pure-Ada TLS 1.3: HTTPS without a C library}\index{TLS}\index{HTTPS}
\label{ch:tls}

The previous chapter put a sockets API over the W5500 and fetched a weather
forecast, over \emph{plain} HTTP, on port 80. That is fine for a hobby query on
a trusted LAN, but the moment the data matters, or the path crosses the open
internet, the answer is HTTPS: TLS over the TCP stream. The W5500 carries TCP/IP
in silicon but knows nothing of TLS, so the entire protocol (the handshake, the
key schedule, certificate validation, the record layer) has to run on our side.

The usual way to get TLS onto a microcontroller is to link a C library: mbedTLS,
BearSSL, wolfSSL. This project does not. The whole point has been to see how far
bare-metal Ada can go on its own, so the TLS client is written in Ada, standing on
two things already in the box: \textbf{SPARKNaCl}\footnote{SPARKNaCl is Roderick
Chapman's SPARK\,2014 reimplementation of TweetNaCl --- itself a compact C edition
of Bernstein, Lange and Schwabe's NaCl. It is proved free of run-time errors in
SPARK, BSD-3-Clause licensed (\textcopyright~2021 Protean Code Limited), and
vendored here under \texttt{crates/sparknacl}. Full credit to that project, without
which this chapter would not exist; see \url{https://github.com/rod-chapman/SPARKNaCl}
and Further Reading (Chapter~\ref{ch:references}).} (a vendored, SPARK-proven Ada
cryptography library) for the public-key and hashing primitives that are
delicate to get right, and the \textbf{ESP32-S3's hardware crypto accelerators}
(AES, SHA, RSA/MPI, the RNG) for the heavy or timing-sensitive arithmetic. What is
left (the protocol itself, the X.509 parser, the P-256 curve) is hand-written
Ada, bounds-checked and profile-portable like everything else.

\cnote{On ESP-IDF you would call \texttt{esp\_tls}, a thin wrapper over mbedTLS,
and never see the handshake. Here there is no library to wrap: the handshake state
machine, the HKDF key schedule, and the certificate chain check are all Ada source
in \texttt{libs/tls}. The cost is that we implement only what we need (TLS~1.3, two
key-exchange groups, two signature families); the benefit is that every byte on
the wire is accounted for in code you can read.}

\section{The shape of the stack}\index{SPARKNaCl}
\label{ch:tls:stack}

TLS is a \emph{record layer} sitting on the TCP stream: after a handshake agrees
on keys, every application byte travels inside an encrypted, authenticated record.
So the TLS client is not a new transport: it is a layer between the application
and \texttt{GNAT.Sockets}, using the same \texttt{Send\_Socket} /
\texttt{Receive\_Socket} from Chapter~\ref{ch:networking} underneath.

\begin{figure}[ht]
\centering
\begin{tikzpicture}[node distance=2mm]
\node[layer, fill=blue!9] (app) {Your application: an HTTPS GET, a JSON reply};
\node[layer, below=of app] (tls) {\texttt{TLS\_Client}: handshake, key schedule, record encrypt/decrypt};
\node[layer, below=of tls] (cert) {\texttt{X509}, \texttt{Cert\_Verify}, \texttt{Chain\_Verify}: certificate validation};
\node[layer, below=of cert] (crypto) {\texttt{P256} (pure Ada) $\;\cdot\;$ SPARKNaCl (X25519, SHA, HKDF) $\;\cdot\;$ HW AES/SHA/RSA/RNG};
\node[layer, below=of crypto] (sock) {\texttt{GNAT.Sockets} over the W5500: the raw TCP byte stream};
\end{tikzpicture}
\caption{The TLS stack. Everything above the socket line is new in this chapter;
the socket line is the networking stack of Chapter~\ref{ch:networking}, unchanged.
The crypto row mixes three sources: a hand-written curve, a vendored SPARK library,
and the chip's accelerators.}
\label{fig:tls-layers}
\end{figure}

The split in that crypto row is deliberate. Table~\ref{tab:tls-crypto} shows who
does what. The rule of thumb: \emph{symmetric} bulk work and big-integer modular
exponentiation go to silicon (fast, and the RSA accelerator is constant-time);
the \emph{asymmetric} curve and protocol logic are Ada because correctness and
auditability matter more there than raw speed, and a one-shot handshake can spend
a few tens of milliseconds on a scalar multiply without anyone noticing.

\begin{table}[ht]
\centering
\small
\begin{tabular}{l l l}
\toprule
\textbf{Primitive} & \textbf{Used for} & \textbf{Where it runs} \\
\midrule
AES-128-GCM        & record encryption        & HW AES + Ada GHASH \\
SHA-256            & transcript, HKDF, HMAC   & SPARKNaCl (HW SHA available too) \\
HKDF / HMAC        & the key schedule         & Ada over SPARKNaCl SHA-256 \\
X25519             & ECDHE key exchange       & SPARKNaCl (constant-time) \\
P-256              & ECDHE + ECDSA verify     & \texttt{P256}: pure Ada \\
RSA ($m^e \bmod n$) & RSA signature recovery  & HW RSA/MPI accelerator \\
RNG                & ephemeral keys, nonces   & HW RNG (entropy source on) \\
\bottomrule
\end{tabular}
\caption{Where each cryptographic primitive runs. There is no ECC accelerator on
the S3, so the P-256 curve is software; everything else has either a SPARK proof
or a hardware engine behind it.}
\label{tab:tls-crypto}
\end{table}

One quiet prerequisite underlies the whole table: \textbf{entropy}. The RNG must be
seeded from a real noise source before it produces key material, or the ephemeral
keys are predictable and the security is theatre. On the S3 that means turning on
the internal entropy source (the 8\,MHz RC oscillator feeding the RNG, plus the SAR
ADC sampling) once at start-up, with \texttt{ESP32S3.RNG.Enable\_Entropy\_Source}.
Every TLS example does this as its first act.

\section{Walking the handshake}\index{TLS handshake}
\label{ch:tls:handshake}

A TLS~1.3 handshake is one round trip. The client opens with a
\texttt{ClientHello}; the server answers with everything at once: its
\texttt{ServerHello}, then (encrypted) its certificate, a signature proving it
holds the certificate's key, and a \texttt{Finished} MAC over the transcript; the
client checks all of that and sends its own \texttt{Finished}; the channel is then
open. \texttt{TLS\_Client.Hello} drives the entire exchange:

\begin{lstlisting}
procedure Hello (S    : in out Session;
                 Sock : GNAT.Sockets.Socket_Type;
                 Host : String;
                 Ok   : out Boolean);
\end{lstlisting}

\subsection{ClientHello: offering what we can do}\index{ClientHello}

The \texttt{ClientHello} advertises exactly the capabilities the client
implements, no more, so the server cannot pick something we cannot honour. We
offer one cipher suite (\texttt{TLS\_AES\_128\_GCM\_SHA256}), TLS~1.3 only, two
key-exchange groups (\texttt{x25519} and \texttt{secp256r1}, with a key share for
each so the server need not ask twice), and the signature schemes we can verify
(\texttt{rsa\_pss\_rsae\_sha256} and \texttt{ecdsa\_secp256r1\_sha256}).

One extension earned its place the hard way. A minimal \texttt{ClientHello} that
omits \texttt{psk\_key\_exchange\_modes} is accepted by lenient servers (a local
Python test server, Cloudflare) but \emph{silently dropped} by others. A real
public server closed the TCP connection right after our hello, with no alert to
explain why. Real clients always send the extension even when offering no
pre-shared key; adding it (\texttt{psk\_dhe\_ke}) was the single change that let
the board talk to the wider internet.

\subsection{The key schedule}\index{HKDF}\index{key schedule}

Once the \texttt{ServerHello} arrives, both sides share an ECDHE secret, and the
TLS~1.3 key schedule (RFC~8446 \S7.1) grinds it through HKDF into a tree of
secrets and keys. The two HKDF operations are \emph{Extract} (an HMAC that mixes a
salt into input keying material) and \emph{Expand-Label} (an HMAC chain that
stretches a secret into named, length-bound output). \texttt{Derive\_Secret}
combines Expand-Label with a hash of the running transcript:

\begin{lstlisting}
Early     := Extract (salt => 0, IKM => 0);
Derived   := Derive_Secret (Early, "derived", Hash (""));
Handshake := Extract (salt => Derived, IKM => Shared_Secret);   --  ECDHE in here
S_HS      := Derive_Secret (Handshake, "s hs traffic", Transcript_Hash);
C_HS      := Derive_Secret (Handshake, "c hs traffic", Transcript_Hash);
\end{lstlisting}

From \texttt{S\_HS} and \texttt{C\_HS} (the server and client handshake traffic
secrets) come the AES keys and IVs that protect the rest of the handshake, and
later the application keys. The schedule is exact and unforgiving: a single wrong
byte anywhere (a mis-ordered transcript, an off-by-one label) yields keys that
do not match, and the server's \texttt{Finished} fails to verify. That property is
also the test: when the board's derived \texttt{S\_HS} matched, byte for byte, the
\texttt{SERVER\_HANDSHAKE\_TRAFFIC\_SECRET} that OpenSSL logged for the same
connection, the X25519 exchange, the transcript hashing, and every label were all
proven correct at once.

\subsection{Two ways to agree on a secret}\index{ECDHE}\index{X25519}\index{P-256}
\label{ch:tls:ecdhe}

The shared secret comes from ECDHE: each side sends an ephemeral public key, and
each multiplies its own private scalar by the other's public point. Modern servers
prefer \textbf{X25519}, which SPARKNaCl provides with a constant-time
implementation, so for that group the client simply calls
\texttt{SPARKNaCl.Scalar.Mult}. But some hardened or FIPS-configured servers
disable X25519 and require the NIST curve \textbf{P-256}, which SPARKNaCl does not
implement. So the project carries its own P-256 in \texttt{libs/tls/p256.adb}:
256-bit modular arithmetic in Montgomery form (the CIOS algorithm), Jacobian point
operations, and a double-and-add scalar multiply. The client offers both groups
and uses whichever the server selects:

\begin{lstlisting}
if S.Group = 16#0017# then                       --  secp256r1 chosen
   Ok := P256.ECDH (S.P256_Priv, Server_X, Server_Y, Shared);
else                                              --  x25519 (the default)
   Shared := SPARKNaCl.Scalar.Mult (S.Priv, S.Server_Pub);
end if;
\end{lstlisting}

The same \texttt{P256} package also verifies ECDSA signatures, which the next
section needs: one curve implementation, two jobs. Its arithmetic is validated
against OpenSSL-generated test vectors by the \texttt{esp32s3\_p256\_kat} example.

\cnote{ECDSA \emph{verification} touches only public values, so the \texttt{P256}
verifier is ordinary variable-time code: there is no secret to leak. ECDH key
exchange does use a secret scalar; the implementation here is still variable-time,
which is acceptable for a fresh ephemeral key used once per handshake but is the
one place a constant-time ladder would be the proper hardening.}

\section{Trusting the certificate}\index{certificate validation}
\label{ch:tls:certs}

Agreeing on keys with a stranger is worthless without knowing \emph{who} the
stranger is. The server proves its identity with a certificate chain and a
signature, and the client must do four things with them: parse the certificate,
check the server actually holds its private key, walk the chain up to a trusted
root, and confirm the certificate is for the host we asked for and is valid right
now.

\subsection{Parsing X.509}\index{X.509}

A certificate is DER-encoded ASN.1: nested tag/length/value triples. The
\texttt{X509} parser is strict by design: it parses attacker-controlled bytes, so
every read is bounds-checked and any malformation yields \texttt{Valid => False}
rather than an out-of-range access. It extracts the signed region (to be hashed),
the validity dates, the subject-alternative-name host list, and the public key,
classifying the key as RSA or P-256 EC, and the signature as one of
\texttt{RSA-SHA256}, \texttt{ECDSA-SHA256}, or \texttt{ECDSA-SHA384}, so the
verifier can dispatch without walking the DER again.

\subsection{Checking the signatures}\index{RSA}\index{ECDSA}\index{Ed25519}\index{RFC 6979}

Two signature families cover essentially all of the public web. \textbf{RSA}
certificates are verified by recovering $m = s^e \bmod n$ on the hardware MPI
accelerator and comparing the recovered padding block: PKCS\#1~v1.5 with
SHA-256, SHA-384 or SHA-512 (an ECC root commonly signs its intermediates with
SHA-384) for chain links, PSS for the TLS \texttt{CertificateVerify}.
\textbf{ECDSA/P-256} and \textbf{Ed25519} certificates are verified in software
(\texttt{P256.Verify}, and SPARKNaCl's Ed25519). The \texttt{CertificateVerify}
step, where the server signs the handshake transcript to prove key possession,
accepts either \texttt{rsa\_pss\_rsae\_sha256} or \texttt{ecdsa\_secp256r1\_sha256};
the chain verifier dispatches each link on how that certificate was signed. The
result is a client that authenticates both an RSA server (Let's Encrypt issues
these) and an ECDSA server (Cloudflare, Google, many CDNs), both confirmed
against the real sites.

The P-256 module also \emph{signs}: \texttt{P256.Sign} produces an ECDSA signature
with a deterministic nonce (RFC~6979, HMAC-SHA-256), so it needs no per-signature
randomness (and the catastrophic nonce-reuse failure mode is gone) while
reproducing the RFC's published vector bit-for-bit. That is the primitive a client
certificate or a TLS \emph{server} side would build on.

\subsection{Anchoring the chain, and knowing the time}\index{certificate chain}\index{trust anchor}\index{key usage}\index{basic constraints}

\texttt{Chain\_Verify.Validate} ties it together: the server's leaf and
intermediates must each be in date, each link's signature must verify under the
next certificate's key, the top must be signed by a \emph{pinned} root the device
trusts, and the leaf's SAN must cover the host. The roots are baked into the
firmware (the HTTPS-weather example pins ISRG~Root~X1, the Let's Encrypt root).
There is no system trust store on a microcontroller.

Anchoring is not the whole of trust. \texttt{Validate} also enforces the
X.509~v3 \emph{usage} extensions a browser checks: every issuing certificate
(each intermediate, and the pinned root when it signs the top) must be a CA
(\texttt{basicConstraints}~\texttt{cA=TRUE}, plus \texttt{keyCertSign} if it
carries a \texttt{keyUsage}) or the chain is rejected as \texttt{Not\_A\_CA}; and
the leaf, if it restricts \texttt{extKeyUsage}, must permit \texttt{serverAuth}
(else \texttt{Bad\_Key\_Usage}). Without the first check a leaf certificate could
pose as an intermediate and sign for any name -- a classic validation hole. The
\texttt{esp32s3\_x509\_chain} self-test drives both rejections with a deliberately
rogue chain (a \texttt{CA:FALSE} certificate that nonetheless signs a leaf), and
the same test exercises the Ed25519 and RSA-SHA-384/512 chains as well.

A parallel rule governs the extensions the parser does \emph{not} recognise.
RFC~5280 requires a certificate be rejected outright if it carries a critical
extension the verifier cannot process --- ``critical'' being the issuer's
declaration that the extension is safety-relevant and \emph{must} be understood, so
silently ignoring one is a validation hole. The X.509 parser therefore records, for
each extension, both whether its \texttt{critical} flag is set (the optional
\texttt{BOOLEAN} between the OID and the value) and whether the parser actually
handled it --- only \texttt{subjectAltName}, \texttt{basicConstraints},
\texttt{keyUsage} and \texttt{extKeyUsage} are recognised. Any extension that is
critical \emph{and} unrecognised sets an \texttt{Unhandled\_Critical} flag on the
certificate record, and that flag is \texttt{and}ed into \texttt{Valid} at the end
of \texttt{Parse}, so an otherwise well-formed certificate carrying a critical
extension the stack does not understand comes back \texttt{Valid = False} and the
chain fails. An unknown \emph{non}-critical extension is still accepted, exactly as
the RFC permits. The \texttt{esp32s3\_x509\_kat} self-test proves the pair with two
otherwise-identical certificates that differ only in that one flag: the critical one
is rejected, the non-critical one accepted.

Validity dates raise a subtlety: ``is this certificate in date?'' needs to know
\emph{today}, and a bare-metal board has no real-time clock. The answer is NTP.
The reusable \texttt{NTP\_Client} (a sibling of \texttt{DNS\_Client}, also pure
\texttt{GNAT.Sockets}) fetches the current UTC time before the handshake; without a
trusted clock the validity check cannot be made, so the client treats an NTP
failure as fatal rather than skipping the check.

\section{The application channel, and a real example}\index{session resumption}\index{NewSessionTicket}\index{AES-GCM}
\label{ch:tls:app}

After the client's \texttt{Finished} is sent, the handshake is complete: fresh
application keys are derived, and \texttt{TLS\_Client.Send} / \texttt{Recv} carry
plaintext in and out, encrypting each record on the way down and decrypting and
authenticating it on the way up. To the application it is just a reliable byte
channel: a GET goes in, a JSON reply comes out.

A server's post-handshake \texttt{NewSessionTicket} is \emph{captured} rather than
skipped, and its resumption key derived, so a later connection can \emph{resume}.
\texttt{TLS\_Client.Resume} runs a second handshake whose ClientHello offers that
ticket as a pre-shared key (still with a fresh key share, so forward secrecy
holds), carrying the PSK binder over the truncated ClientHello. If the server
accepts, the handshake is shorter and sends no certificate flight at all: there is
nothing to re-validate, only a \texttt{Finished} to check. The
\texttt{esp32s3\_tls\_resume} example shows the two connections back to back, and
is verified against \texttt{openssl s\_server}.

The \texttt{esp32s3\_tls\_weather} example is the whole chapter in one program. It
fetches a live forecast from \texttt{api.open-meteo.com} \emph{over HTTPS}:

\begin{lstlisting}
[wx] resolving api.open-meteo.com ...           -- DNS_Client (UDP)
[wx] api.open-meteo.com = 188.40.99.226
[wx] getting time from NTP ...                  -- NTP_Client (UDP)
[wx] NTP UTC = 2026-06-26 03:33:58
[wx] TLS 1.3 up: cipher 0x1301                  -- ECDHE + key schedule
[wx] CertificateVerify (RSA-PSS): OK            -- server proved key possession
[wx] server Finished: OK
[wx] chain validation to ISRG Root X1: 2 certs -> VALID
[wx] forecast for 52.52, 13.41 (HTTPS):
[wx]   temperature : 19.2 C
[wx]   wind speed  : 5.7 km/h
\end{lstlisting}

Every layer of this chapter is in those lines: DNS and NTP over the UDP path, the
TLS~1.3 handshake with its ECDHE and HKDF schedule, an RSA-PSS
\texttt{CertificateVerify} against the site's real Let's Encrypt certificate, a
chain validated to a pinned root at an NTP-supplied time, and finally an encrypted
GET whose decrypted JSON is scraped for the temperature, all in Ada, with no C
TLS library linked, on a board whose Ethernet chip has never heard of TLS.

\section{What is implemented, and what is not}
\label{ch:tls:status}

Honesty about scope matters in a security chapter. The client implements TLS~1.3
with the \texttt{TLS\_AES\_128\_GCM\_SHA256} suite, X25519 and P-256 key exchange,
and certificate authentication across the algorithms real chains use: RSA
(PKCS\#1 and PSS) with SHA-256/384/512, ECDSA-P256 with SHA-256/384, and Ed25519.
Chain validation is full: signatures, validity at an NTP-supplied time, the leaf
hostname, the basic-constraints / key-usage / extended-key-usage policy checks,
and anchoring to pinned roots. The stack also \emph{signs} (deterministic
ECDSA-P256) and \emph{resumes} a session from a PSK ticket. That is enough to
fetch from the large majority of real public HTTPS servers, as the open-meteo and
Cloudflare runs show.

It deliberately stops short of a few things a production stack would add: the
\texttt{TLS\_AES\_256\_GCM\_SHA384} and ChaCha20-Poly1305 cipher suites (the first
needs SHA-384 wired through the key schedule, the second a ChaCha20-Poly1305 that
SPARKNaCl does not provide); certificate revocation (OCSP/CRL); a TLS \emph{server}
side; and client certificates (the ECDSA signing primitive is in place, but the
handshake wiring is not). None of these are hard in principle (each is a bounded
slice on top of the same primitives), but each is a place where ``it compiled''
is not the same as ``it is safe,'' which is exactly why they are called out here
rather than quietly assumed.
